\begin{proposition}[2.5.8]
\label{IV.2.5.8-fr}
Sean $A$ un anillo conmutativo semi-local, $B$ una $A$-álgebra (no
necesariamente conmutativa), y $V$ y $W$ dos $B$-módulos. Sea $A'$ una
$A$-álgebra conmutativa que es un $A$-módulo \textup{fielmente plano}; se pone $B'=B\otimes_A A'$, $V'=V\otimes_A A'$, $W'=W\otimes_A A'$, de modo que $B'$ es una $A'$-álgebra, y $V'$ y $W'$ son $B'$-módulos.
Se supone además que se cumple una de las condiciones siguientes:
\begin{enumerate}
\item[\textup{(i)}] $A$ y $A'$ son noetherianos, y $V$ y $W$ son $A$-módulos de
tipo finito.
\item[\textup{(ii)}] $B$ es un $A$-módulo de tipo finito, y $V$ es un $B$-módulo
proyectivo de tipo finito y un $A$-módulo de presentación finita.
\end{enumerate}
Entonces, si $V'$ y $W'$ son $B'$-módulos isomorfos, $V$ y $W$ son
$B$-módulos isomorfos.
\end{proposition}

Observemos que, en el caso (ii), como $W'$ es $A'$-isomorfo a $V'$, es un
$A'$-módulo de tipo finito, de donde resulta que $W$ es un $A$-módulo de tipo
finito (Bourbaki, \emph{Alg. comm.}, cap.~I, \S~3,
núm.~6, prop.~11); por tanto, en todos los casos, $V$ y $W$
son $A$-módulos de tipo finito. Además:

\begin{env}[2.5.8.1]
\label{IV.2.5.8.1-fr}
\itshape Bajo una de las hipótesis \textup{(i), (ii)}, $\operatorname{Hom}_B(V,W)$ es un
$A$-módulo de tipo finito.
\end{env}

Es evidente en el caso (i), pues $\operatorname{Hom}_A(V,W)$ es entonces un
$A$-módulo de tipo finito y $\operatorname{Hom}_B(V,W)$ es un sub-$A$-módulo
de $\operatorname{Hom}_A(V,W)$. En el caso (ii), $V$ es sumando directo de un
$B$-módulo libre $B_s^n$, por lo que $\operatorname{Hom}_B(V,W)$ es sumando
directo de $\operatorname{Hom}_B(B_s^n,W)=W^n$, y, como $W$ es un $A$-módulo de
tipo finito, lo mismo ocurre con $\operatorname{Hom}_B(V,W)$. Pongamos
\[
C=\operatorname{End}_B(V),\qquad M=\operatorname{Hom}_B(V,W)
\]
que son $A$-módulos de tipo finito en los casos (i) y (ii). Se sabe que, bajo
una de las condiciones (i), (ii), el homomorfismo canónico

\begin{env}[2.5.8.2]
\label{IV.2.5.8.2-fr}
\[
\operatorname{Hom}_A(V,W)\otimes_A A'\longrightarrow
\operatorname{Hom}_{A'}(V',W')
\]
\end{env}

es \emph{biyectivo} (Bourbaki, \emph{Alg. comm.}, cap.~II, \S~2,
núm.~10, prop.~11). Como $A'$ es un $A$-módulo plano,
$\operatorname{Hom}_B(V,W)\otimes_A A'$ se identifica canónicamente con un
sub-$A'$-módulo de $\operatorname{Hom}_A(V,W)\otimes_A A'$. La imagen de este
submódulo por el homomorfismo (2.5.8.2) está contenida en
$\operatorname{Hom}_{B'}(V',W')$, pues, si $u\in\operatorname{Hom}_B(V,W)$ y
$a'\in A'$, la imagen de $u\otimes a'$ por (2.5.8.2) es el homomorfismo
$u':V'\to W'$ tal que $u'(x\otimes1)=u(x)\otimes a'$; para todo $b\in B$,
se tiene, por tanto,
$u'((b\otimes1)(x\otimes1))=u'(bx\otimes1)=u(bx)\otimes a'=bu(x)\otimes a'=(b\otimes1)(u(x)\otimes a')$,
de donde se sigue nuestra afirmación. Dicho esto:

\begin{env}[2.5.8.3]
\label{IV.2.5.8.3-fr}
\itshape Bajo una de las hipótesis \textup{(i), (ii)}, el homomorfismo
\end{env}

\begin{env}[2.5.8.4]
\label{IV.2.5.8.4-fr}
\[
\operatorname{Hom}_B(V,W)\otimes_A A'\longrightarrow
\operatorname{Hom}_{B'}(V',W')
\]
\itshape es biyectivo.
\end{env}

Para todo $b\in B$, denotemos por $h(b)$ (resp. $h'(b)$) la homotecia
$x\leadsto bx$ en $V$ (resp. $W$), que es un $A$-endomorfismo. Sea
$(b_\alpha)_{\alpha\in I}$ un sistema de generadores de la $A$-álgebra $B$;
la aplicación
\[
u\leadsto(h'(b_\alpha)\circ u\setmin u\circ h(b_\alpha))_{\alpha\in I}
\]
de $\operatorname{Hom}_A(V,W)$ en $(\operatorname{Hom}_A(V,W))^I$ es
$A$-lineal y, por definición, su \emph{núcleo} no es sino
$\operatorname{Hom}_B(V,W)$; en otras palabras, se tiene una sucesión exacta
\[
0\longrightarrow\operatorname{Hom}_B(V,W)
\longrightarrow\operatorname{Hom}_A(V,W)
\longrightarrow(\operatorname{Hom}_A(V,W))^I.
\]
El mismo razonamiento se aplica reemplazando $A$, $B$, $V$, $W$ por $A'$,
$B'$, $V'$, $W'$; además, se tiene un diagrama

\begin{env}[2.5.8.5]
\label{IV.2.5.8.5-fr}
\[
\xymatrix{
0 \ar[r] & \operatorname{Hom}_B(V,W)\otimes_A A' \ar[r] \ar[d]_{r} &
\operatorname{Hom}_A(V,W)\otimes_A A' \ar[r] \ar[d]_{s} &
(\operatorname{Hom}_A(V,W))^I\otimes_A A' \ar[d]_{t}\\
0 \ar[r] & \operatorname{Hom}_{B'}(V',W') \ar[r] &
\operatorname{Hom}_{A'}(V',W') \ar[r] &
(\operatorname{Hom}_{A'}(V',W'))^I
}
\]
\end{env}

donde $r$ es el homomorfismo (2.5.8.4), $s$ el homomorfismo (2.5.8.2) y
$t$ el homomorfismo compuesto
\[
(\operatorname{Hom}_A(V,W))^I\otimes_A A'
\xrightarrow{w}(\operatorname{Hom}_A(V,W)\otimes_A A')^I
\xrightarrow{s^I}(\operatorname{Hom}_{A'}(V',W'))^I
\]
siendo $w$ el homomorfismo canónico (Bourbaki, \emph{Alg.}, cap.~II,
3.ª ed., \S~3, núm.~7). Se comprueba de
inmediato que el diagrama (2.5.8.5) es conmutativo y, como $A'$ es un
$A$-módulo plano, sus filas son exactas. Por último, se ha visto

\oldpage[IV-2]{26}
que $s$ es un isomorfismo, y lo mismo ocurre, por tanto, con $s^I$; en el caso
(ii), se puede tomar $I$ \emph{finito}, y se sabe entonces que $w$ es biyectivo
(Bourbaki, \emph{loc. cit.}, prop.~7); en el caso (i), se observa que, si $B'$
(resp. $B''$) es la imagen de $B$ por $h$ (resp. $h'$) en
$\operatorname{End}_A(V)$ (resp. $\operatorname{End}_A(W)$), $B'$ y $B''$
son $A$-módulos de tipo finito, por lo que también se puede tomar $I$ \emph{fini}to.
Así, en todos los casos, $t$ es biyectivo, y se concluye, por tanto, que lo
mismo ocurre con $r$.

Por consiguiente, de (2.5.8.4) resulta que, si se pone
\[
C'=C\otimes_A A',\qquad M'=M\otimes_A A'
\]
se tienen biyecciones canónicas
